\section*{第 \ref{ch:g10:earth-energy-budget} 章——地球的能量收支}

\begin{solution}{exo:g10:earth-energy-budget:1}
某地的\omterm{def:g10:earth-energy-budget:insolation}{日射量}是指其上空大气层顶部、单位水平面积所接收的
太阳光功率（按一日或一年平均）。某表面的\omterm{def:g10:earth-energy-budget:albedo}{反照率}是指落在其上的
太阳光中被其反射的比例。新雪（0.8 至 0.9）反射
的比例远高于开阔海面
（0.06）。
\end{solution}

\begin{solution}{exo:g10:earth-energy-budget:2}
反射：$100/340 \approx \qty{29}{\percent}$；被大气吸收：
$80/340 \approx \qty{24}{\percent}$；被地表吸收：$160/340
\approx \qty{47}{\percent}$。
\end{solution}

\begin{solution}{exo:g10:earth-energy-budget:3}
沥青：$800 \times (1 - 0.04) = \qty{768}{W}$。雪：$800 \times (1 -
0.85) = \qty{120}{W}$，少了六倍以上。
\end{solution}

\begin{solution}{exo:g10:earth-energy-budget:4}
蒸发为 \qty{82}{W/m^2}，远大于空气增温的
\qty{21}{W/m^2}。蒸发的水从海洋中取走\omterm{def:g10:earth-energy-budget:latent-heat}{潜热}；水汽随
空气上升，在往往高出数千米的地方凝结成云滴，
并将这部分能量交还给周围的空气。
\end{solution}

\begin{solution}{exo:g10:earth-energy-budget:5}
十二月北半球倾向远离太阳：在极点
太阳全天位于地平线以下，因此没有阳光到达。十二月
冬至日，这种情况在 $90^\circ - 23.44^\circ =
\qty{66.56}{\degree}$ 北纬以北的整个地区都是如此。
\end{solution}

\begin{solution}{exo:g10:earth-energy-budget:6}
$1361 \sin 90^\circ = \qty{1361}{W}$；$1361 \sin 30^\circ =
\qty{680.5}{W}$；$1361 \sin 10^\circ \approx \qty{236}{W}$。$S_0$ 的一半
则需 $\sin h = 0.5$，即 $h = \qty{30}{\degree}$。
\end{solution}

\begin{solution}{exo:g10:earth-energy-budget:7}
反射 100，被大气吸收 80，被地表吸收 160：
$100 + 80 + 160 = 340$，阳光的去向全部算清。在地表，
$185 = 25 + 160$：未被反射的部分即被吸收。
\end{solution}

\begin{solution}{exo:g10:earth-energy-budget:8}
圆盘面积：$\pi \times (6.371\times10^6)^2 \approx \qty{1.275e14}{m^2}$；功率：
$1361 \times 1.275\times10^{14} \approx \qty{1.74e17}{W}$。球面：$4\pi
R^2 \approx \qty{5.10e14}{m^2}$。平均值：$1.74\times10^{17}/5.10\times
10^{14} \approx \qty{340}{W/m^2}$，约为 $S_0$ 的四分之一。
\end{solution}

\begin{solution}{exo:g10:earth-energy-budget:9}
$0.7 \times 5.10\times10^{14} \times 3.16\times10^7 \approx
\qty{1.1e22}{J}$，为一年内所得。
\end{solution}

\begin{solution}{exo:g10:earth-energy-budget:10}
云层多反射了 $100 - 53 = \qty{47}{W/m^2}$ 的阳光，这部分
因而未被吸收。云还阻挡了 $267 - 239 = \qty{28}{W/m^2}$ 的
红外辐射。综合而言，地球失去 $47 - 28 = \qty{19}{W/m^2}$：云层
总体上使地球降温。
\end{solution}

\begin{solution}{exo:g10:earth-energy-budget:11}
约自 5 月 10 日至 7 月底，即 5 月、6 月和 7 月。在
赤道，太阳在两个分点日正午位于天顶，即 3 月和
9 月；在至日时，太阳偏离天顶 \qty{23.44}{\degree}，
因此赤道在 6 月和 12 月接收到的太阳能略少。
\end{solution}

\begin{solution}{exo:g10:earth-energy-budget:12}
每秒：$82/(2.47\times10^6) \approx \qty{3.32e-5}{kg}$。每天：
$3.32\times10^{-5} \times 86\,400 \approx \qty{2.9}{kg}$，形成约
\qty{2.9}{mm} 厚的水层。每年：$3.32\times10^{-5} \times 3.16\times10^7 \approx
\qty{1050}{kg}$，约合 \qty{1.05}{m} 深的水。
\end{solution}

\begin{solution}{exo:g10:earth-energy-budget:13}
由 $S_0 \sin\delta = S_0/\pi$ 得 $\sin\delta = 1/\pi \approx 0.318$，因此
$\delta \approx \qty{18.6}{\degree}$。然后 $\sin\lambda = 0.318/\sin 23.44^\circ
\approx 0.318/0.398 \approx 0.800$：$\lambda \approx \qty{53.15}{\degree}$ 或
$180^\circ - 53.15^\circ = \qty{126.85}{\degree}$。两者之间，$\lambda$
增加了 \qty{73.7}{\degree}，即一年的 $73.7/360$，约
\num{75} 天，即从 5 月中旬到 7 月下旬。
\end{solution}

\begin{solution}{exo:g10:earth-energy-budget:14}
冰吸收了 $1 - 0.6 = 0.40$ 的阳光，海洋吸收 $1 - 0.06 =
0.94$。额外功率：$(0.94 - 0.40) \times 300 \times 10^{12} \approx
\qty{1.6e14}{W}$，超过全球植物光合作用的总功率。额外的
能量使海水变暖，融化了浮冰边缘的更多冰，
使更多暗色海水裸露，这一变化会自我加强。
\end{solution}

\begin{solution}{exo:g10:earth-energy-budget:15}
最近时：$1361 \times (149.6/147.1)^2 \approx \qty{1408}{W/m^2}$；最远时：
$1361 \times (149.6/152.1)^2 \approx \qty{1317}{W/m^2}$；相差约
\qty{7}{\percent}。在北纬 \qty{60}{\degree} 处，六月的\omterm{def:g10:earth-energy-budget:insolation}{日射量}
是十二月的二十倍：几个百分点的变化无法
造成这样的差异。此外，地球在一月初离太阳最近，正值
北半球冬季；而两个半球与太阳的距离相同，
却季节相反：造成季节的是地轴的倾斜，
而不是距离。
\end{solution}

\begin{solution}{pb:g10:earth-energy-budget:1}
\textbf{1.} $\pi \times (6.371\times10^6)^2 \approx \qty{1.275e14}{m^2}$。

\textbf{2.} $1361 \times 1.275\times10^{14} \approx \qty{1.735e17}{W}$.

\textbf{3.} $4\pi R^2 \approx \qty{5.101e14}{m^2}$.

\textbf{4.} $1.735\times10^{17}/5.101\times10^{14} \approx \qty{340}{W/m^2}$,
即 $S_0/4$。

\textbf{5.} 赤道：\qty{1361}{W}；北纬 \qty{60}{\degree}：$1361 \sin
30^\circ = \qty{680.5}{W}$。

\textbf{6.} $100/340 \approx 0.294$。

\textbf{7.} $340 - 100 = \qty{240}{W/m^2}$；对整个地球，$240 \times
5.101\times10^{14} \approx \qty{1.224e17}{W}$。

\textbf{8.} $53/340 \approx 0.156$。云使\omterm{def:g10:earth-energy-budget:albedo}{反照率}增大的倍数为
$100/53 \approx 1.9$：它们几乎使其翻倍。

\textbf{9.} $(240/5.67\times10^{-8})^{1/4} \approx \qty{255}{K}$，约
比地表的 \qty{288}{K} 低 \qty{33}{K}。

\textbf{10.} $160 + 342 = \qty{502}{W/m^2}$。

\textbf{11.} $82 + 21 + 398 = \qty{501}{W/m^2}$。地表获得的能量约为
\qty{1}{W/m^2}：即四舍五入后的 \qty{0.7}{W/m^2} 失衡量。

\textbf{12.} $342/160 \approx 2.1$。大气传给地表的红外能量
是地表所吸收太阳光能量的两倍多：
这就是温室效应。

\textbf{13.} $82 \times 3.156\times10^7/(2.47\times10^6) \approx
\qty{1050}{kg}$，相当于厚约 \qty{1.05}{m} 的一层。

\textbf{14.} $0.7 \times 5.101\times10^{14} \approx \qty{3.6e14}{W}$。

\textbf{15.} $2805/72.07 \approx \qty{38.9}{kJ}$，即每克碳。

\textbf{16.} $104.9\times10^9$ 吨即 $1.049\times10^{17}$~g；
$1.049\times10^{17} \times 38.9\times10^3 \approx \qty{4.08e21}{J}$。

\textbf{17.} $4.08\times10^{21}/3.156\times10^7 \approx \qty{1.29e14}{W}$。

\textbf{18.} $1.29\times10^{14}/5.101\times10^{14} \approx
\qty{0.25}{W/m^2}$。

\textbf{19.} $0.7/0.25 \approx 2.8$：地球的失衡量几乎是
光合作用所储存功率的三倍。

\textbf{20.} $1.29\times10^{14}/1.224\times10^{17} \approx 0.0011$：
光合作用储存的能量仅占地球吸收阳光的\textbf{约 \qty{0.1}{\percent}}
——即千分之一。
\end{solution}
